\documentclass{article}
\usepackage[utf8]{inputenc}
\usepackage[T1]{fontenc}
\usepackage{amsmath,amssymb,bm}
\usepackage[margin=2cm]{geometry}
\usepackage{graphicx}
\usepackage{subcaption}
\captionsetup{font=small, labelfont=bf}
\usepackage{float}
\usepackage[section]{placeins}
\usepackage[numbers,square,sort&compress]{natbib}
\usepackage{hyperref}
\graphicspath{{figures/}}
\renewcommand{\topfraction}{0.9}
\renewcommand{\bottomfraction}{0.8}
\renewcommand{\textfraction}{0.07}
\renewcommand{\floatpagefraction}{0.85}

\title{Event Cameras in Robotics: An Overview}
\author{Amit Nativ}
\date{\today}

\begin{document}

\maketitle

\section{Event cameras}

Event cameras are bio-inspired vision sensors that do not capture frames at a fixed rate.
Instead, each pixel independently reports changes in brightness as a stream of asynchronous
\emph{events}. This gives them microsecond temporal resolution, a high dynamic range
(140\,dB vs.\ 60\,dB for frame cameras), low power consumption and little motion
blur~\cite{gallegoEventbasedVisionSurvey2022,chakravarthiRecentEventCamera2024}.

\paragraph{Event generation model.}
Each pixel responds to the log intensity $L=\log I$. An event $e_k=(\mathbf{x}_k,t_k,p_k)$ is
triggered at pixel $\mathbf{x}_k=(x_k,y_k)^\top$ and time $t_k$ as soon as the brightness
change since the last event at that pixel reaches a contrast threshold
$C>0$~\cite{gallegoEventbasedVisionSurvey2022}:
\begin{equation}
    \Delta L(\mathbf{x}_k,t_k) \doteq L(\mathbf{x}_k,t_k) - L(\mathbf{x}_k,t_k-\Delta t_k) = p_k\,C ,
    \label{eq:event_generation}
\end{equation}
where $\Delta t_k$ is the time since the previous event at $\mathbf{x}_k$ and the polarity
$p_k\in\{+1,-1\}$ is the sign of the change (Figure~\ref{fig:generation}(a)). For small
$\Delta t_k$, a first-order Taylor expansion shows that each event measures the temporal
derivative of brightness, $\partial L/\partial t\,(\mathbf{x}_k,t_k)\approx p_kC/\Delta t_k$.
Under constant illumination and brightness constancy, the increment becomes
\begin{equation}
    \Delta L \approx -\nabla L\cdot\mathbf{v}\,\Delta t ,
    \label{eq:moving_edges}
\end{equation}
where $\mathbf{v}$ is the image-plane velocity. Events are therefore caused by moving edges: a
pixel fires only where there is texture ($\nabla L\neq\mathbf{0}$) and the motion has a
component normal to the edge ($\nabla L\cdot\mathbf{v}\neq 0$). The output is a sparse,
asynchronous point set in space--time rather than an image
(Figure~\ref{fig:generation}(b)), so frame-based vision algorithms do not apply directly.
Figure~\ref{fig:davis} shows a real event camera, the DAVIS, and its output.


\begin{figure}[!htb]
    \centering
    \includegraphics[width=\textwidth]{event_generation.pdf}
    \caption{\textbf{Event generation.} (a)~Log intensity $L(\mathbf{x}_0,t)$ at one pixel
    and the events it triggers. An ON (OFF) event fires each time $L$ rises (falls) by $C$
    from its level at the last event (dashed), as in~\eqref{eq:event_generation}; gray lines
    mark the levels $L(\mathbf{x}_0,0)+kC$. A 30\,fps camera samples the same signal only
    three times (squares). (b)~Events of a bright dot rotating at 5 revolutions per second,
    plotted in space--time: the dot traces a helix, with ON events on its leading edge and OFF
    events on its trailing edge. Both panels are simulated with the noise-free
    model~\eqref{eq:event_generation}.}
    \label{fig:generation}
\end{figure}

\begin{figure}[!htb]
    \centering
    \includegraphics[width=0.6\textwidth]{paper_gallego_davis.png}
    \caption{\textbf{The DAVIS event camera.} (a)~Simplified circuit of a DAVIS pixel, which
    combines an event pixel (DVS, red) with a conventional frame pixel (APS, blue).
    (b)~DVS operation: ON and OFF events fire when the log intensity crosses the threshold
    around its reset level. (c)--(d)~The DAVIS chip and USB camera. (e)~A white square on a
    rotating black disk produces a spiral of events in space--time (green: past, red:
    present), while frames sample it at a fixed rate. (f)~A frame overlaid with 5\,ms of
    events. Reproduced from~\cite{gallegoEventbasedVisionSurvey2022}.}
    \label{fig:davis}
\end{figure}

\section{Processing events}
\label{sec:processing}

Most methods first aggregate a window of events $\mathcal{E}=\{e_k\}_{k=1}^{N_e}$ into a grid
that conventional or spiking networks can process~\cite{gallegoEventbasedVisionSurvey2022}.
Figure~\ref{fig:representations} shows three common choices. The \emph{event frame} sums the
polarities at each pixel; scaled by $C$, it approximates the brightness increment over the
window, by~\eqref{eq:event_generation}:
\begin{equation}
    \Delta\hat{L}(\mathbf{x}) = C\sum_{k=1}^{N_e} p_k\,\delta(\mathbf{x}-\mathbf{x}_k) .
    \label{eq:event_frame}
\end{equation}
The \emph{time surface} keeps only the most recent timestamp $t_{\mathrm{last}}(\mathbf{x})$ at
each pixel, weighted by an exponential kernel so that recent motion dominates:
\begin{equation}
    \mathcal{T}(\mathbf{x},t) = \exp\!\left(-\frac{t-t_{\mathrm{last}}(\mathbf{x})}{\tau}\right) .
    \label{eq:time_surface}
\end{equation}
The \emph{voxel grid} keeps more of the timing. It splits the window into $B$ temporal bins
and spreads each event's polarity over the two nearest bins:
\begin{equation}
    V(\mathbf{x},b) = \sum_{k=1}^{N_e} p_k\,\delta(\mathbf{x}-\mathbf{x}_k)\,
    \max\bigl(0,\,1-|b-t_k^*|\bigr),
    \qquad t_k^* = (B-1)\,\frac{t_k-t_1}{t_{N_e}-t_1},
    \label{eq:voxel_grid}
\end{equation}
with $b\in\{0,\dots,B-1\}$. E-RAFT, for example, uses voxel grids as the input to a dense
optical-flow network~\cite{gehrigERAFTDenseOptical2021}. Figure~\ref{fig:paper_representations}
shows these representations on real data.

\begin{figure}[!htb]
    \centering
    \includegraphics[width=\textwidth]{event_representations.pdf}
    \caption{\textbf{Grid representations of the same 250\,ms of events}, simulated for a
    scene translating at $\mathbf{v}=(60,25)$\,px/s. (a)~Event frame~\eqref{eq:event_frame}:
    edges smear along their trajectories. (b)~Time surface~\eqref{eq:time_surface} with
    $\tau=50$\,ms: recent events are dark, so the fading trail shows the direction of motion.
    (c)~Voxel grid~\eqref{eq:voxel_grid} with $B=5$: bins 0, 2 and 4 show the edges at three
    times. Red: positive polarity; blue: negative.}
    \label{fig:representations}
\end{figure}

\begin{figure}[!htb]
    \centering
    \includegraphics[width=\textwidth]{paper_gallego_representations.png}
    \caption{\textbf{Event representations of real data} (\emph{slider\_depth} sequence).
    (a)~Events in space--time; here positive events are blue and negative events red, the
    opposite of Figures~\ref{fig:representations} and~\ref{fig:lif}. (b)~Event frame.
    (c)~Time surface of the negative events (dark: recent). (d)~Interpolated voxel grid with
    10 bins. (e)~Motion-compensated image of warped events
    (Section~\ref{sec:processing}, \emph{Contrast maximization}). (f)~Intensity image
    reconstructed from the events. Reproduced from~\cite{gallegoEventbasedVisionSurvey2022}.}
    \label{fig:paper_representations}
\end{figure}

\paragraph{Contrast maximization.}
An event frame blurs because each edge fires along its whole trajectory.
\citet{gallegoUnifyingContrastMaximization2018} turn this into an estimation principle: warp
each event to a reference time $t_{\mathrm{ref}}$ along a candidate motion $\bm{\theta}$, and
choose the motion that makes the resulting image of warped events (IWE) sharpest. For optical
flow, $\bm{\theta}=\mathbf{v}$ and
\begin{align}
    \mathbf{x}'_k &= \mathbf{x}_k - (t_k-t_{\mathrm{ref}})\,\bm{\theta} ,
    \label{eq:warp}\\
    H(\mathbf{x};\bm{\theta}) &= \sum_{k=1}^{N_e} b_k\,\delta(\mathbf{x}-\mathbf{x}'_k) ,
    \label{eq:iwe}
\end{align}
where $b_k=1$ counts events and $b_k=p_k$ sums polarities. Sharpness is measured by the
variance of the IWE over its $N_p$ pixels $h_{ij}$, with mean $\mu_H$:
\begin{equation}
    f(\bm{\theta}) = \sigma^2\bigl(H(\mathbf{x};\bm{\theta})\bigr)
    = \frac{1}{N_p}\sum_{i,j}\bigl(h_{ij}-\mu_H\bigr)^2 ,
    \qquad
    \bm{\theta}^* = \arg\max_{\bm{\theta}} f(\bm{\theta}) .
    \label{eq:cmax}
\end{equation}
Figure~\ref{fig:cmax} shows the objective for the events of Figure~\ref{fig:representations}.
With other warps, the same objective estimates depth and rotational
motion~\cite{gallegoUnifyingContrastMaximization2018}. It needs no ground truth, and it is the
self-supervised loss with which a small drone learns depth on
board~\cite{hagenaarsOnDeviceSelfSupervisedLearning2025}.

\begin{figure}[!htb]
    \centering
    \includegraphics[width=\textwidth]{contrast_maximization.pdf}
    \caption{\textbf{Contrast maximization for optical flow} on the events of
    Figure~\ref{fig:representations} ($b_k=1$, bilinear voting). (a)~Without warping
    ($\bm{\theta}=\mathbf{0}$), events pile up along their trajectories and the IWE is blurred.
    (b)~Warping with the optimal flow aligns them into sharp edges. (c)~The
    variance~\eqref{eq:cmax} over candidate flows has a single peak at
    $\bm{\theta}^*=(60,25)$\,px/s, equal to the true flow.}
    \label{fig:cmax}
\end{figure}

\section{Event cameras in robotics}

In robotics, the main advantage of event cameras is low latency.
\citet{falangaHowFastToo2019} model a robot flying at constant speed $v$ toward an obstacle that
it detects at range $s$, with perception latency $t_{\mathrm{lat}}$. To clear the obstacle, the
robot must move sideways by $d$ with lateral acceleration bounded by $a_{\max}$; the
time-optimal (bang--bang) manoeuvre takes $2\sqrt{d/a_{\max}}$. Avoidance is safe only if the
time left after detection covers this manoeuvre:
\begin{equation}
    \frac{s}{v} - t_{\mathrm{lat}} \;\geq\; 2\sqrt{\frac{d}{a_{\max}}}
    \qquad\Longrightarrow\qquad
    v_{\max} = \frac{s}{t_{\mathrm{lat}} + 2\sqrt{d/a_{\max}}} .
    \label{eq:vmax}
\end{equation}
Latency therefore directly lowers the maximum safe speed, and it matters most for agile
platforms, whose manoeuvre time is short.

Event cameras have enabled the first quadrotor flight with event-based state estimation, which
fused events, frames and inertial measurements to fly in low light and high-dynamic-range
scenes~\cite{vidalUltimateSLAMCombining2018}. They also enabled position control of a quadrotor
yawing at over 20\,rad/s after a rotor failure, where events outperformed frames in low
light~\cite{sunAutonomousQuadrotorFlight2021}. In reactive tasks, a quadrotor avoided
fast-moving obstacles at relative speeds of up to 10\,m/s with 3.5\,ms
latency~\cite{falangaDynamicObstacleAvoidance2020} (Figure~\ref{fig:falanga}), an FPGA-based collision-avoidance pipeline
reached 2.14\,ms end to end~\cite{bonazziLowLatencyEventbasedObstacle2025a}, and a quadruped
caught balls thrown at up to 15\,m/s~\cite{forraiEventbasedAgileObject2023}. Combining events
with a 20\,fps camera gives the latency of a 5{,}000\,fps camera at the bandwidth of a 45\,fps
one~\cite{gehrigLowlatencyAutomotiveVision2024}.

\begin{figure}[!htb]
    \centering
    \includegraphics[width=0.85\textwidth]{paper_falanga_avoidance.png}
    \caption{\textbf{Avoiding a dynamic obstacle with an event camera.} A quadrotor detects a
    ball thrown at it with its onboard event camera and dodges it. Reproduced
    from~\cite{falangaDynamicObstacleAvoidance2020}.}
    \label{fig:falanga}
\end{figure}

\section{Learning and neuromorphic processing}

Recent work learns event-based policies end to end. Quadrotors avoid static obstacles using
only a monocular event camera~\cite{bhattacharyaMonocularEventBasedVision2024}
(Figure~\ref{fig:bhattacharya}) and fly through
clutter at up to 9.8\,m/s~\cite{messikommerApproximateImitationLearning}, and a small drone
learns monocular depth from events on board~\cite{hagenaarsOnDeviceSelfSupervisedLearning2025}.
Training data is a recurring obstacle. Simulated events are expensive to render and lack a
faithful continuous-time sensor model, so these works rely on approximate simulated events,
converting ordinary video into events~\cite{gehrigVideoEventsRecycling2019}, fine-tuning on real
data, or self-supervised learning.

\begin{figure}[!htb]
    \centering
    \includegraphics[width=0.62\textwidth]{paper_bhattacharya_stream.png}
    \caption{\textbf{Learning obstacle avoidance from a monocular event camera.} Part of the
    continuous event stream of a forest flight (5\% of the events shown) is split into 33\,ms
    batches; the network maps each binary event mask to a depth prediction $D(\theta)$ and an
    avoidance command. Reproduced from~\cite{bhattacharyaMonocularEventBasedVision2024}.}
    \label{fig:bhattacharya}
\end{figure}

Events are also a natural input for spiking neural networks (SNNs), which communicate with
binary spikes. Counting the events of a patch in each time step gives the input $X[t]$ of a
leaky integrate-and-fire (LIF) neuron~\cite{eshraghianTrainingSpikingNeural2023}:
\begin{equation}
    U[t] = \beta U[t-1] + W X[t] - S_{out}[t-1]\,\Theta ,
    \qquad
    S_{out}[t] =
    \begin{cases}
        1, & U[t] \geq \Theta\\
        0, & \text{otherwise}
    \end{cases}
    \label{eq:lif_events}
\end{equation}
The neuron therefore spikes only while events arrive and stays silent otherwise
(Figure~\ref{fig:lif}). On neuromorphic processors, such networks run event by event. An SNN
on a neuromorphic chip has controlled a UAV~\cite{vitaleEventdrivenVisionControl2021}, and a
fully neuromorphic network on Intel's Loihi mapped raw events to low-level control commands for
a freely flying drone at 200\,Hz and 27\,$\mu$J per
inference~\cite{paredes-vallesFullyNeuromorphicVision2023}. Its spiking vision network is
trained self-supervised with contrast maximization~\eqref{eq:cmax} on real events
(Figure~\ref{fig:paredes}). Processing each event as it
arrives, without grouping events into frames, has been demonstrated for mobile-robot
dodging~\cite{jiangFullyAsynchronousNeuromorphic2024}.

\begin{figure}[!htb]
    \centering
    \includegraphics[width=\textwidth]{lif_events.pdf}
    \caption{\textbf{LIF neuron driven by events.} (a)~Events from a $7\times7$-pixel patch
    as the bright disk of Figure~\ref{fig:representations} passes it: ON events at its leading
    edge, OFF events at its trailing edge. (b)~Membrane potential~\eqref{eq:lif_events}, with
    $X[t]$ the number of events in each 2\,ms step, $W=0.05$, $\beta=0.9$ and $\Theta=1$.
    (c)~Output spikes: the neuron fires only while events arrive.}
    \label{fig:lif}
\end{figure}

\begin{figure}[!htb]
    \centering
    \includegraphics[width=0.85\textwidth]{paper_paredes_snn.png}
    \caption{\textbf{Spiking vision network of a fully neuromorphic drone.} Running at about
    200\,Hz on the neuromorphic processor, accumulated events from each image corner pass
    through three spiking encoders and a pooling layer, which output the optical flow of that
    corner. During training, the four corner flows give a dense flow field, which deblurs the
    events for the self-supervised contrast-maximization loss. Reproduced
    from~\cite{paredes-vallesFullyNeuromorphicVision2023}.}
    \label{fig:paredes}
\end{figure}

\bibliographystyle{unsrtnat}
\bibliography{references}

\end{document}
